\documentclass[11pt,a4paper]{article}
\usepackage{lmodern}
\usepackage[T1]{fontenc}
\usepackage[utf8]{inputenc}
\usepackage[margin=2.2cm]{geometry}
\usepackage{listings}
\usepackage{xcolor}
\usepackage{needspace}
\usepackage[colorlinks=true,linkcolor=black,urlcolor=blue]{hyperref}
\usepackage[expansion=false]{microtype}

\lstset{
  basicstyle=\ttfamily\footnotesize,
  breaklines=true,
  breakatwhitespace=false,
  columns=fullflexible,
  keepspaces=true,
  frame=single,
  rulecolor=\color{black!30},
  backgroundcolor=\color{black!3},
  aboveskip=6pt,
  belowskip=6pt,
  xleftmargin=2pt,
  xrightmargin=2pt
}

\setlength{\parindent}{0pt}
\setlength{\parskip}{6pt}

\begin{document}

\begin{center}
{\Large\bfseries CS3423 Compilers 2\par}
\vspace{4pt}
{\large\bfseries Mini Assignment 3 -- LLVM IR Exploration Report\par}
\end{center}

\noindent Names: Krishnan R, Harikrishna\\
Roll Numbers: CS24BTECH11036, CS24BTECH11030\\
Platform: WSL2 (Ubuntu)\\
Clang Version: Clang version 17.0.6\\
LLVM Version: 17.0.6 (built from source: llvmorg-17.0.6 with clang and clang-tools-extra, X86 target)

\tableofcontents

\section*{Introduction}

This report delves into the LLVM compilation pipeline using Clang 17. It involves generating an IR (Intermediate Representation) for various constructs, analysing control flow, data structures and floating-point values (all at -O0).


\section*{Methodology}

LLVM and Clang 17.0.6 were built from source (llvmorg-17.0.6 on release/17.x) with only the clang and clang-tools-extra projects enabled, in Release mode, for the X86 target. Every .c/.cpp file was compiled to LLVM IR using clang -S -emit-llvm -O0 -o file.ll file.c (clang++ for C++), with optimizations disabled so that the IR maps directly to the source. The x86-64 assembly of the floating-point program was generated using llc -O0. All the outputs can be regenerated by running run.sh from the submission directory. The folders follow the assignment items: Q2 (IR study), Q3 (data structures) and Q4 (floating point). Item 1 is theory only, so it has no code.

\section{LLVM Directory Layout}

The LLVM source (release/17.x of llvm-project) is divided into the following main directories:


\textbf{llvm/: }contains the core of LLVM. include/llvm holds the header files, lib/ holds their implementations (IR, Analysis, Transforms (optimization passes), CodeGen, Target), tools/ holds programs like llc, opt and lli, and test/ holds the regression tests.


\textbf{clang/: }contains the C-family frontend. lib/Lex, lib/Parse, lib/Sema and lib/AST handle lexing, parsing, semantic analysis and the AST, while lib/CodeGen converts the AST to LLVM IR (the .ll files in this report).


\textbf{clang-tools-extra/: }contains extra tools built on Clang such as clang-tidy and clangd.


\textbf{Other directories }(cmake/, runtimes/, ...) hold build files and runtime libraries. Only clang and clang-tools-extra are enabled using -DLLVM\_ENABLE\_PROJECTS="clang-tools-extra;clang".


Flow followed in this report: C/C++ source $\rightarrow$ clang frontend $\rightarrow$ LLVM IR (.ll) $\rightarrow$ llc $\rightarrow$ assembly.


\section{LLVM IR Study of Core C/C++ Constructs}

Four C programs were used: Q2\_variables.c (variables), Q2\_conditionals\_loops.c (conditionals and loops), Q2\_functions.c (functions) and Q2\_casts.c (data types and casts).

\subsection{C Code in Q2\_variables.c (variables)}

\begin{lstlisting}
#include <stdio.h>
int a = 5; // (i) Global Integer
int main(){
        int b = 10; // (ii) A Local Integer
        int c = 3;  // (ii) A Local Integer
        int d = a + b * c; // (iii) A local Integer initialized to a + b * c
}
\end{lstlisting}

\medskip\noindent\textbf{IR Code Snippets in Q2\_variables.ll:}\par

\begin{lstlisting}
@a = dso_local global i32 5, align 4 ; (i) Global Integer a = 5
  %1 = alloca i32, align 4           ; (ii) allocate space for Local Integer b
  %2 = alloca i32, align 4           ; (ii) allocate space for Local Integer c
  %3 = alloca i32, align 4
  store i32 10, ptr %1, align 4      ; (ii) Initialize a Local Integer b = 10
  store i32 3, ptr %2, align 4       ; (ii) Initialize a Local Integer c = 3
  %4 = load i32, ptr @a, align 4     ; load a
  %5 = load i32, ptr %1, align 4     ; load b
  %6 = load i32, ptr %2, align 4     ; load c
  %7 = mul nsw i32 %5, %6            ; b * c
  %8 = add nsw i32 %4, %7            ; a + b * c
  store i32 %8, ptr %3, align 4      ; (iii) Store a + b * c in d
\end{lstlisting}

\textbf{Representation of local variables}


In LLVM IR, named registers follow the SSA (Static Single Assignment) form, meaning each register is assigned exactly once and does not have a memory address. On the other hand, C variables are real variables stored in memory (stack) which can be done using alloca (which allocates memory on stack to the variable), have an address, can be modified multiple times throughout the program. Therefore, Clang represents local variables using alloca instructions, which allocate space on the stack, and uses load and store instructions to read and modify their values. This allows LLVM IR to correctly model the mutable nature of variables in C. 


\subsection{C Code in Q2\_conditionals\_loops.c (conditionals and loops)}

\begin{lstlisting}
#include <stdio.h>
int main(){
        int a = 2, b = 3;
        // (i) If-Else Condition
        if (a == 2){
                a = 2;
        }
        else{
                a = 2;
        }
        // (ii) Switch Case Condition
        switch(a){
                case 1: a = 2; break;
                case 2: a = 2; break;
                default: a = 2; break;
        }
        // (iii) For loop
        for (int i = a; i < 5; i++){
                a++;
        }
        // (iv) While loop
        while (a < 5){
                a++;
        }
        // (v) Do-While loop
        do {
                a++;
        } while ( a < 5);
        return 0;
}
\end{lstlisting}

\medskip\noindent\textbf{IR Code Snippets in Q2\_conditionals\_loops.ll:}\par

The IR of each construct is shown below. For the loops, the same loop was written in all three forms (for, while and do-while) and their IR was compared.


\begin{lstlisting}
Preds = %x1 %x2... means that the given block can be reached from blocks x1, x2, ...
\end{lstlisting}

\textbf{Explanation of each construct:}


\textbf{if-else}


\begin{lstlisting}
  %5 = load i32, ptr %2, align 4   ; load a
  %6 = icmp eq i32 %5, 2          ; check: a == 2
  br i1 %6, label %7, label %8    ; if (true) -> go to block 7 else -> go to block 8
7:                                ; preds = %0, Label 7 starts here
  store i32 2, ptr %2, align 4    ; do: a = 2 then
  br label %9                     ; go to next part (Block 9)
8:                                ; preds = %0, Label 8 starts here
  store i32 2, ptr %2, align 4    ; do: a = 2 then
  br label %9                     ; go to next part (Block 9)
\end{lstlisting}

\textbf{switch}

\begin{lstlisting}
9:                                ; preds = %8, %7, Label 9 starts here
  %10 = load i32, ptr %2, align 4 ; load variable a into ptr 10
  switch i32 %10, label %13 [     ; switch(ptr 10 = a): else -> go to 13
    i32 1, label %11              ; if a == 1 -> go to 11
    i32 2, label %12              ; else if a == 2 -> go to 12
  ]
11:                               ; preds = %9, Block 11 starts here
  store i32 2, ptr %2, align 4    ; Set a = 2
  br label %14                    ; Go to block 14
12:                               ; preds = %9, Label 12 starts here
  store i32 2, ptr %2, align 4    ; Set a = 2
  br label %14                    ; Go to block 14
13:                               ; preds = %9, Label 13 starts here
  store i32 2, ptr %2, align 4    ; Set a = 2
  br label %14                    ; Go to block 14
\end{lstlisting}

\textbf{for:}


\begin{lstlisting}
14:                                 ; preds = %13, %12, %11, Label 14 starts here
  %15 = load i32, ptr %2, align 4   ; read value from %2 (variable a), store it in temporary %15
  store i32 %15, ptr %4, align 4    ; store that value into %4
  br label %16                      ; Jump to loop start
16:                                 ; preds = %22, %14, Label 16 starts here
  %17 = load i32, ptr %4, align 4   ; Load variable i into %17
  %18 = icmp slt i32 %17, 5         ; Check if i < 5
  br i1 %18, label %19, label %25 ; If i < 5, branch to loop body, else branch to exit loop
19:                                 ; preds = %16, Label 19, loop body starts here
  %20 = load i32, ptr %2, align 4   ; Load variable a into %20
  %21 = add nsw i32 %20, 1          ; Store a++ into %21
  store i32 %21, ptr %2, align 4    ; Store the new value of a
  br label %22                      ; Branch to label 22
22:                                 ; preds = %19, Label 22 starts here
  %23 = load i32, ptr %4, align 4   ; Load variable i into %23
  %24 = add nsw i32 %23, 1          ; Store i++ into %24
  store i32 %24, ptr %4, align 4    ; Store the new value of i
  br label %16                      ; Branch to label 16
25:                                 ; preds = %16, Label 25 starts here
  br label %26                      ; Branch to label 26
\end{lstlisting}

\textbf{while}


\begin{lstlisting}
26:                               ; preds = %29, %25, Label 26, while loop starts here
  %27 = load i32, ptr %2, align 4 ; Load the value of variable a into %27
  %28 = icmp slt i32 %27, 5       ; Check if a < 5
  br i1 %28, label %29, label %32 ; Go to loop runner if a < 5, else go to exit loop
29:                               ; preds = %26, Label 29, Loop runner starts here
  %30 = load i32, ptr %2, align 4 ; Load the value of variable a into %30
  %31 = add nsw i32 %30, 1        ; a++
  store i32 %31, ptr %2, align 4  ; Store the new value of a
  br label %26                    ; Branch to label 26 for the next iteration
32:                               ; preds = %26, Label 32, Loop exit starts here
  br label %33                    ; Branch to Label 33
\end{lstlisting}

\textbf{do-while}


\begin{lstlisting}
33:                        ; preds = %36, %32, Label 33, Do-While loop starts here
  %34 = load i32, ptr %2, align 4 ; Load the value of variable a into %34
  %35 = add nsw i32 %34, 1        ; a++
  store i32 %35, ptr %2, align 4  ; Store the new value of a
  br label %36                    ; Branch to label 36
36:                               ; preds = %33, Label 36 starts here
  %37 = load i32, ptr %2, align 4 ; Load the value of variable a into %37
  %38 = icmp slt i32 %37, 5       ; Check if a < 5
  br i1 %38, label %33, label %39 ; If a < 5, repeat the loop, else exit
39:                               ; preds = %36, Label 39, loop exit starts here
  ret i32 0                       ; int main() returns 0
\end{lstlisting}

\textbf{Comparison of the for and while loops}


The for and while loops are almost identical in their respective IR's. Both are implemented using a conditional (icmp) and branch (br) block, a loop body, branch instruction to either repeat the loop or exit the loop. The only difference is when we write the for loop in C, we explicitly define the loop iterator variable i, increment steps in the syntax while in the while loop, the compiler takes care of these details.


\textbf{Position of the condition check in the do-while loop}


In both the for and while loops the condition check appears before entering the loop body while in the do-while loop they appear after the loop body finishes one iteration.


\subsection{C Code in Q2\_functions.c (functions)}

\begin{lstlisting}
#include <stdio.h>
int square (int x){
        return x * x;
}
int main(){
        printf("The square of 3 = %d", square(3));
        return 0;
}
\end{lstlisting}

\medskip\noindent\textbf{IR Code Snippets in Q2\_functions.ll:}\par




\textbf{define and declare}


The define keyword is present as define dso\_local i32 @square(i32 noundef \%0) \#0 (The definition of the int square(int) function) and as define dso\_local i32 @main() \#0  (The definition of the main() function)


The declare keyword is present as declare i32 @printf(ptr noundef, ...) \#1 (Tells the compiler the function printf() definition is not present in the current file, since it is present in the external library stdio.h included as a header in the file)


The define keyword is used to define functions defined in the given file (like int square (int x) and int main()) itself while the declare keyword is used to declare functions whose definitions are not present in the current file but rather in external files (like printf() in the external library stdio.h declared at the start of the program) but called from the given file. 


\textbf{Passing arguments by value and returning the value}




\begin{lstlisting}
define dso_local i32 @square(i32 noundef %0) #0 {
  %2 = alloca i32, align 4         ; Allocate memory to store the argument in %0
  store i32 %0, ptr %2, align 4    ; Receives argument x and stores it in %2
  %3 = load i32, ptr %2, align 4
  %4 = load i32, ptr %2, align 4
  %5 = mul nsw i32 %3, %4
  ret i32 %5         ; Return the value after performing the required calculations
}
%2 = call i32 @square(i32 noundef 3) ; We pass the value 3 into the square() function by value. The call instruction is used to receive the return value conveyed back by the function.
\end{lstlisting}

When main() calls square(3), the value 3 is passed directly to the square() function and received as the parameter \%0. It is stored in \%2 (allocated using alloca), loaded twice and multiplied to get \%5, which is returned using ret i32 \%5. In main(), the call instruction gives this returned value in register \%2, which is then passed to printf (\%3 = call i32 (ptr, ...) @printf(ptr noundef @.str, i32 noundef \%2)).


\textbf{The ret instruction}




\begin{lstlisting}
ret i32 %5   ; return result of square
ret i32 0    ; return from main
\end{lstlisting}

The ret function is used to terminate a function, return the control to the caller while also returning value specified. In the square function ret i32 \%5 returns the value x * x while in the main function ret i32 0 returns the value 0 to the OS.


\subsection{C Code in Q2\_casts.c (data types and casts)}

\begin{lstlisting}
#include <stdio.h>
int main(){
        // (i) Assigning Double to an Int (Narrowing)
        int a = 1;
        double b = 2.0;
        a = (int) b;
        // (ii) Assigning Int to a Long (Widening)
        int c = 3;
        long d = 4;
        d = (long) c;
        // (iii) Assigning Long to an Int (Narrowing)
        c = (int) d;
        return 0;
}
\end{lstlisting}

\medskip\noindent\textbf{IR Code Snippets in Q2\_casts.ll:}\par

The cast instructions used (fptosi, sext, trunc) and when each one is chosen:


\textbf{fptosi}




\begin{lstlisting}
%6 = load double, ptr %3, align 8
%7 = fptosi double %6 to i32
\end{lstlisting}

The instruction fptosi (floating point to signed integer) is used to convert a double to an integer (narrowing) where precision is lost in the process of converting a larger type (double) to a smaller type (int).


\needspace{6\baselineskip}
\textbf{sext}




\begin{lstlisting}
%8 = load i32, ptr %4, align 4
%9 = sext i32 %8 to i64
\end{lstlisting}

The instruction sext (sign extension) is used to convert a 32-bit integer to a 64-bit integer (widening) where the sign is preserved while the value is extended.


\textbf{trunc}




\begin{lstlisting}
  %10 = load i64, ptr %5, align 8
  %11 = trunc i64 %10 to i32
\end{lstlisting}

The instruction trunc (truncate) is used to convert a larger integer type (i64, long) to a smaller one (i32, int) by discarding the higher bits, so the value can change if it does not fit.


\section{Lowering of Data Structures to LLVM IR}

3(a) IR was generated for arrays (3.1), std::vector (3.2), structs (3.3), unions (3.4) and classes (3.5). The questions 3(b), 3(c) and 3(d) are answered in 3.6, 3.7 and 3.8.

\subsection{C Code in Q3\_array.c (arrays)}

\begin{lstlisting}
#include <stdio.h>
#include <stdlib.h>
int main(){
        int arr1[5];
        int *arr2 = (int*) malloc(5 * sizeof(int));
        for(int i = 0; i < 5; i++){
                arr1[i] = i;
        }
        int d = arr1[2];
        for(int i = 0; i < 5; i++){
                arr2[i] = i;
        }
        int e = arr2[2];
        return 0;
}
\end{lstlisting}

\medskip\noindent\textbf{IR Code Snippets in Q3\_array.ll:}\par




\textbf{Array access using getelementptr (GEP)}




The getelementptr (GEP) instruction is used to compute the address of an element in an array. GEP does not perform memory access, it only computes addresses. 


Ex: \%16 = getelementptr inbounds [5 x i32], ptr \%2, i64 0, i64 \%15


\begin{lstlisting}
[5 x i32]  Array type (5 integers)
ptr %2  Base address of array
i64 0  First index (array itself)
i64 %15  Element index (i)
arr1[i] -> address = base + i * sizeof(int)  GEP Computes this address
\end{lstlisting}

\textbf{Allocation of an array (alloca)}




The alloca instruction looks like \%2 = alloca [5 x i32], align 16 for an array type.


The alloca instruction allocates memory for the array on the stack. Here, the [5 x i32] indicates that the memory allocated is equal to the memory size of 5 integers.


On the other hand, the scalar alloca looks like \%1 = alloca i32, align 4 where the memory allocated is equal to the memory size of 1 integer.


\textbf{Static and dynamic (malloc) allocation}




Static Allocation: \%2 = alloca [5 x i32], align 16


Memory equal to the memory size of 5 integers is statically allocated on the stack for the array arr1.


\begin{lstlisting}
declare noalias ptr @malloc(i64 noundef) #1
%8 = call noalias ptr @malloc(i64 noundef 20)
store ptr %8, ptr %3
\end{lstlisting}

Since malloc() is an external function, a declare instruction is used.


Memory equal to the memory size of 5 integers (5 integers * 4 bytes = 20 bytes) is dynamically allocated on the heap for the array arr2 using malloc which returns a pointer \%8 whose value is then stored in ptr \%3.


LLVM IR represents array access using the getelementptr instruction, which computes addresses based on indices. Static arrays are allocated on the stack using alloca, while dynamic arrays are allocated on the heap using malloc. This difference is reflected in the IR through distinct allocation mechanisms and pointer types used in element access. 


\subsection{C++ Code in Q3\_vector.cpp (std::vector)}

\begin{lstlisting}
#include <vector>
using namespace std;
int main(){
        vector<int> v;
        v.push_back(1);
        v.push_back(2);
        return v[0];
}
\end{lstlisting}

\medskip\noindent\textbf{IR Code Snippets in Q3\_vector.ll (attributes shortened to ...):}\par\nopagebreak[4]

\begin{lstlisting}
%"class.std::vector" = type { %"struct.std::_Vector_base" }
%"struct.std::_Vector_base<int, std::allocator<int>>::_Vector_impl_data" = type { ptr, ptr, ptr }  ; begin, end, end_of_storage
%2 = alloca %"class.std::vector", align 8             ; vector object (3 pointers)
call void @_ZNSt6vectorIiSaIiEEC2Ev(ptr ... %2)         ; constructor
invoke void @_ZNSt6vectorIiSaIiEE9push_backEOi(ptr ... %2, ptr ... %3) to label %7 unwind label %12   ; v.push_back(1)
%9 = call ptr @_ZNSt6vectorIiSaIiEEixEm(ptr ... %2, i64 noundef 0)   ; v[0]: pointer to element
%10 = load i32, ptr %9, align 4
call void @_ZNSt6vectorIiSaIiEED2Ev(ptr ... %2)         ; destructor
; body of operator[]:
%8 = load ptr, ptr %7, align 8                           ; load the begin pointer
%10 = getelementptr inbounds i32, ptr %8, i64 %9         ; &begin[index]
\end{lstlisting}

std::vector is not a built-in type in LLVM IR. It is a normal C++ class, so it is represented as a struct (\%"class.std::vector") whose innermost type is three pointers (begin, end and end\_of\_storage). The vector object itself takes only 24 bytes on the stack, while the elements are stored on the heap. Operations like the constructor, push\_back, operator[] and the destructor are all calls to functions, and operator[] just loads the begin pointer and uses a getelementptr with one index to get the address of the element. The push\_back calls use invoke (instead of call) with a landingpad since they can throw an exception. The line using namespace std; only changes how the names are looked up by the compiler, so it adds no IR.


\subsection{C Code in Q3\_struct.c (struct)}

\begin{lstlisting}
#include <stdio.h>
struct Point{
        int x;
        int y;
};
int main(){
        struct Point p;
        p.x = 1;
        p.y = 2;
        return 0;
}
\end{lstlisting}

\medskip\noindent\textbf{IR Code Snippets in Q3\_struct.ll:}\par\nopagebreak[4]

\textbf{Type definition and member alignment}

Type definition of struct Point\{ int x; int y; \}: \%struct.Point = type \{ i32, i32 \}

The type definition shows that the struct Point consists of two i32 members (x and y) stored sequentially in memory.

\begin{lstlisting}
%2 = alloca %struct.Point, align 4
store i32 1, ptr %3, align 4
\end{lstlisting}

Member alignment is not specified in the type definition but is specified during memory allocation (alloca), and memory access (store, load) using (align) keyword. Each i32 is aligned to 4 bytes.

\textbf{Accessing a struct member vs an array element}

\begin{lstlisting}
Struct access:
%3 = getelementptr inbounds %struct.Point, ptr %2, i32 0, i32 0
%4 = getelementptr inbounds %struct.Point, ptr %2, i32 0, i32 1
\end{lstlisting}

Array access: \%15 = getelementptr inbounds [5 x i32], ptr \%2, i64 0, i64 \%index

In Struct access, we can access only fixed field indices (compile time constants). The GEP type is \%struct.Point while in Array access we can access variable field indices (i, runtime values). The GEP type is [5 x i32].

\subsection{C Code in Q3\_union.c (union)}

\begin{lstlisting}
union Data{
        int i;
        float f;
};
int main(){
        union Data d;
        d.i = 1;
        d.f = 1.5f;
        return 0;
}
\end{lstlisting}

\medskip\noindent\textbf{IR Code Snippets in Q3\_union.ll:}\par

\begin{lstlisting}
%union.Data = type { i32 }                  ; union type: only one i32
%2 = alloca %union.Data, align 4            ; size 4, alignment 4
store i32 1, ptr %2, align 4                ; d.i = 1
store float 1.500000e+00, ptr %2, align 4   ; d.f = 1.5f  (same address, different type)
\end{lstlisting}

LLVM IR has no separate union type. A union is represented as a struct with one member (here i32). All the members are accessed using the same address \%2, only the type of the load/store changes (i32 for d.i and float for d.f).


\subsection{C++ Code in Q3\_class.cpp (class)}

\begin{lstlisting}
class Point{
public:
        int x;
        int y;
        // Constructor
        Point(){
                x = 0;
                y = 0;
        }
};
int main(){
        Point p;
        return 0;
}
\end{lstlisting}

\medskip\noindent\textbf{IR Code Snippets in Q3\_class.ll (attributes shortened):}\par\nopagebreak[4]

\textbf{Type definition and the constructor}

Type definition of class Point\{ int x; int y; \}: \%class.Point = type \{ i32, i32 \}

\begin{lstlisting}
call void @_ZN5PointC2Ev(ptr %2)
define void @_ZN5PointC2Ev(ptr %0)
\end{lstlisting}

Constructors appear as normal functions in LLVM IR. The object pointer (this: \%0) is passed as an argument, and the constructor initializes the fields using getelementptr and store instructions. Initializing happens with a constructor call with address of p (Here \%2 in the call function).

\subsection{Arrays vs vectors}

\textbf{3(b) What is the difference between arrays and vectors in LLVM IR?}

An array is a built-in type of LLVM IR ([5 x i32]) whose size is fixed at compile time and is allocated directly using alloca. Its elements are accessed using getelementptr with two indices (0, i). A vector (std::vector) is a class, so it is a struct of three pointers (24 bytes) and its elements are stored on the heap. It is accessed only through function calls (push\_back, operator[]), where operator[] uses a getelementptr with one index on the begin pointer. Note: LLVM IR also has its own vector type (\textless{}4 x i32\textgreater{}) used for SIMD, which is different from std::vector.

\subsection{Type of a union}

\textbf{3(c) How is the type of a union determined in LLVM IR?}

Clang takes the member with the largest alignment (the largest size if there is a tie) as the type of the union and adds padding at the end if the union is bigger than that member. Here i and f both need 4 bytes with alignment 4, so the first one (i32) is used and no padding is needed, giving \{ i32 \} (size 4, alignment 4). If a bigger member such as char c[8] is added, the type becomes \{ i32, [4 x i8] \} (size 8, alignment 4), because c[8] has alignment 1 but needs 8 bytes.

\subsection{struct vs class}

\textbf{3(d) What are the differences between struct and class of C++ once lowered in LLVM IR?}

Q3\_struct\_vs\_class.cpp puts the struct from Q3\_struct.c and the class from Q3\_class.cpp in one file. They are renamed StructPoint and ClassPoint, since a struct and a class cannot both be called Point in the same file.

\begin{lstlisting}
#include <stdio.h>
// struct from Q3_struct.c
struct StructPoint{
        int x;
        int y;
};
// class from Q3_class.cpp
class ClassPoint{
public:
        int x;
        int y;
        // Constructor
        ClassPoint(){
                x = 0;
                y = 0;
        }
};
int main(){
        struct StructPoint sp;
        sp.x = 1;
        sp.y = 2;
        ClassPoint cp;
        return 0;
}
\end{lstlisting}

\medskip\noindent\textbf{IR Code Snippets in Q3\_struct\_vs\_class.ll (attributes shortened to ...):}\par\nopagebreak[4]

\begin{lstlisting}
%struct.StructPoint = type { i32, i32 }
%class.ClassPoint = type { i32, i32 }
%2 = alloca %struct.StructPoint, align 4
%3 = alloca %class.ClassPoint, align 4
%4 = getelementptr inbounds %struct.StructPoint, ptr %2, i32 0, i32 0   ; sp.x = 1
store i32 1, ptr %4, align 4
%5 = getelementptr inbounds %struct.StructPoint, ptr %2, i32 0, i32 1   ; sp.y = 2
store i32 2, ptr %5, align 4
call void @_ZN10ClassPointC2Ev(ptr ... %3)                              ; ClassPoint cp
; body of the constructor:
%4 = getelementptr inbounds %class.ClassPoint, ptr %3, i32 0, i32 0
store i32 0, ptr %4, align 4
%5 = getelementptr inbounds %class.ClassPoint, ptr %3, i32 0, i32 1
store i32 0, ptr %5, align 4
\end{lstlisting}

Both types have the same layout \{ i32, i32 \}. The only difference is the name of the type (\%struct.StructPoint and \%class.ClassPoint). The constructor of ClassPoint is a normal function (\_ZN10ClassPointC2Ev) which takes the pointer to the object as its first argument, while the members of the struct are set directly using getelementptr and store.

There are no differences in the IR once lowered, apart from the name of the type (\%struct.StructPoint and \%class.ClassPoint). Public/private access is checked by the compiler frontend and is not present in the IR. Both are laid out in the same way and their member functions (like the constructor) are normal functions taking the object pointer as an argument. Differences only come from the contents of the type (constructors, virtual functions, base classes) and not from the keyword.

\section{Floating-Point Representation in LLVM IR}

\subsection{C Code in Q4\_float.c}

\begin{lstlisting}
int main(){
        float a = 2.3f;
        double b = 5.7;
        b = b + a;               // fpext + fadd
        a = (float) b * a;       // fptrunc + fmul
        int c = (int) (b - a);   // fpext + fsub + fptosi
        if(b < a){               // fpext + fcmp
                double d = c / a;    // sitofp + fdiv + fpext
                a = -a;              // fneg
        }
        return (int) c;
}
\end{lstlisting}

\medskip\noindent\textbf{IR Code Snippets in Q4\_float.ll:}\par

\begin{lstlisting}
%1 = alloca i32, align 4                       ; return value of main
%2 = alloca float, align 4                     ; float a
%3 = alloca double, align 8                    ; double b
%4 = alloca i32, align 4                       ; int c
%5 = alloca double, align 8                    ; double d
store i32 0, ptr %1, align 4                   ; return slot = 0
store float 0x4002666660000000, ptr %2, align 4 ; a = 2.3f
store double 5.700000e+00, ptr %3, align 8     ; b = 5.7

; b = b + a
%6 = load double, ptr %3, align 8              ; load b
%7 = load float, ptr %2, align 4               ; load a
%8 = fpext float %7 to double                  ; a: float -> double
%9 = fadd double %6, %8                        ; b + a
store double %9, ptr %3, align 8               ; b = b + a

; a = (float) b * a
%10 = load double, ptr %3, align 8             ; load b
%11 = fptrunc double %10 to float              ; (float) b: double -> float
%12 = load float, ptr %2, align 4              ; load a
%13 = fmul float %11, %12                      ; (float) b * a
store float %13, ptr %2, align 4               ; a = result

; int c = (int) (b - a)
%14 = load double, ptr %3, align 8             ; load b
%15 = load float, ptr %2, align 4              ; load a
%16 = fpext float %15 to double                ; a: float -> double
%17 = fsub double %14, %16                     ; b - a
%18 = fptosi double %17 to i32                 ; (int): double -> int
store i32 %18, ptr %4, align 4                 ; c = result

; if (b < a)
%19 = load double, ptr %3, align 8             ; load b
%20 = load float, ptr %2, align 4              ; load a
%21 = fpext float %20 to double                ; a: float -> double
%22 = fcmp olt double %19, %21                 ; b < a (ordered, less than)
br i1 %22, label %23, label %31                ; true -> block 23, false -> 31

23:                                            ; body of the if
; double d = c / a
%24 = load i32, ptr %4, align 4                ; load c
%25 = sitofp i32 %24 to float                  ; c: int -> float
%26 = load float, ptr %2, align 4              ; load a
%27 = fdiv float %25, %26                      ; c / a
%28 = fpext float %27 to double                ; result: float -> double
store double %28, ptr %5, align 8              ; d = result

; a = -a
%29 = load float, ptr %2, align 4              ; load a
%30 = fneg float %29                           ; -a
store float %30, ptr %2, align 4               ; a = -a
br label %31                                   ; leave the if

31:                                            ; after the if
; return (int) c
%32 = load i32, ptr %4, align 4                ; load c (already int, no cast)
ret i32 %32                                    ; return c
\end{lstlisting}

\subsection{Representation of floating-point numbers}

LLVM IR has float (32 bit), double (64 bit) and others like half, x86\_fp80 and fp128. C float and double map to float and double directly. A constant is written in decimal (5.700000e+00) if the decimal text converts back to exactly the stored value, otherwise as a hexadecimal value (2.3f is written as 0x4002666660000000, the bit pattern of the double holding the nearest float to 2.3).


\subsection{Operations on floating-point values}

The operations supported are fadd, fsub, fmul, fdiv, frem and fneg, and both operands must have the same type. Comparison is done using fcmp which returns an i1 used by br. The predicates are ordered (oeq, olt, ogt, ...) which are false if an operand is NaN, and unordered (une, ult, ...) which are true if an operand is NaN. Clang uses olt for \textless{} and une for !=.


\subsection{Type lowering of floating-point values}

LLVM IR does not do implicit conversions, so Clang adds an explicit instruction for each one. In b + a, the float is first converted to double using fpext, (float) b uses fptrunc, double to int uses fptosi and int to float (in c / a) uses sitofp. When lowered further using llc (Q4\_float.s), these become x86-64 SSE instructions: addsd/subsd for double and mulss/divss for float, cvtss2sd/cvtsd2ss for fpext/fptrunc, cvttsd2si/cvtsi2ss for fptosi/sitofp, ucomisd for fcmp, and an xor with the sign bit (0x80000000) for fneg.


\section*{Declaration of usage of AI/LLMs}

AI tools were used for converting our explanations to LaTeX/PDF, wording changes, checking whether our programs cover the required constructs, verifying results, understanding LLVM IR and clearing doubts. All code was written by us. LLVM was installed and all outputs were generated and checked by us. AI tools were not used to copy.

\section*{References}

[1] LLVM Project, "LLVM Language Reference Manual", \url{https://llvm.org/docs/LangRef.html}


[2] LLVM Project, "Getting Started with the LLVM System" (directory layout, building from source), \url{https://llvm.org/docs/GettingStarted.html}


[3] LLVM Project, "The Often Misunderstood GEP Instruction", \url{https://llvm.org/docs/GetElementPtr.html}


\end{document}
